% ============================================================
% Auto-generated from docs/golden-sunflowers/ch-12-hardware-bridge-deferred.md
% Source of truth: Railway phd-postgres-ssot ssot.chapters (gHashTag/trios#380)
% ============================================================

\chapter{Hardware Bridge (deferred)}

% Chapter Anchor header (Phase 1 UNIFY task 1.4 · trios#380)
% Trinity S^3AI strand · 35 chapters running parallel to the Flos Aureus petals
\begin{tcolorbox}[colback=blue!3,colframe=blue!40!black,title={\textbf{Trinity S\textsuperscript{3}AI Strand} \textbf{Ch.12}}]
  \textbf{Strand:} Trinity S\textsuperscript{3}AI --- silicon, software, science \\
  \textbf{Anchor:} \(\varphi^{2} + \varphi^{-2} = 3\) (Trinity Identity, INV-22) \\
  \textbf{Lane:} S12 (Trinity strand) \\
  \textbf{Theorems in chapter:} 0 \\
  \textbf{Coq link:} \filepath{trinity-clara/proofs/igla/} (per-theorem) \\
  \textbf{Notation key:} GF(16) ternary algebra, IGLA training stack, ASHA pruning; INV-k via \citetheorem{INV-k} (AP.F)
\end{tcolorbox}

\begin{figure}[H]
\centering
\makebox[\linewidth][c]{\includegraphics[width=1.18\linewidth,keepaspectratio]{\figChTwelveHardwareBridge}}
\caption*{Figure — Ch.12: Hardware Bridge (deferred).}
\end{figure}

\begin{quote}\itshape
``Make each program do one thing well.  To do a new job, build afresh rather than
complicate old programs by adding new features.  Expect the output of every
program to become the input to another.''
\upshape --- Doug McIlroy, \textit{Unix Philosophy} (1978)
\end{quote}

\section*{Where the theorem meets the wire}

At some point, a formally verified arithmetic system has to stop being a document
and start being a circuit.  For Trinity S\textsuperscript{3}AI, that moment
happens at a row of 16-bit AXI-Lite registers on a QMTech XC7A100T board---
92~MHz clock, sub-watt power budget, 0 DSP slices.  The Hardware Bridge is the
two-hundred-line RTL layer that stands between the software proofs and the
physical pins.  It is deliberately thin.  Thin bridges do not break.

The bridge's three-channel structure is not arbitrary.  The GoldenFloat exponent
field splits naturally into sub-unity (\(\varphi^{-2}\)), unity, and super-unity
(\(\varphi^2\)) bands---three lanes mandated by the same anchor identity
\(\varphi^2 + \varphi^{-2} = 3\) that organises the weight lattice.  One 16-bit
data channel per band means the host can route token batches to the correct
hardware lane without any format conversion overhead.  The arithmetic structure
visible in the Coq proofs is therefore also visible on the PCB.  This is
co-design in the literal sense: the mathematics and the hardware are the same
decision, expressed in two different languages.

The chapter is marked ``deferred'' because two of its three formal obligations
require FPGA measurements that post-date the mathematical chapters.
Ch.28 will supply the synthesis report (63~toks/sec at 92~MHz) and Ch.31 the
system-level integration tests.  What the present chapter supplies is the
interface contract: signal names, timing assumptions, UART-V6 token-transfer
protocol, and clock-domain crossing rules.  Think of it as the interface
specification that makes it possible to write the proofs before the board
arrives---the intellectual equivalent of programming to an API rather than to an
implementation.

The rest of this chapter defines the bridge architecture (Section~2), the AXI-Lite
register map (Section~3), the UART-V6 protocol (Section~4), error handling and
clock-domain crossing (Section~5), and the forward references to Ch.28 and Ch.31
that complete the verification chain (Section~6).  The DARPA 3000\texttimes{}
energy ratio becomes measurable only once this interface is locked down; this
chapter locks it down.

\section{Abstract}\label{ch_12:abstract}

The Hardware Bridge chapter specifies the interface layer between the Trinity S³AI software stack and the QMTech XC7A100T FPGA. It defines the AXI-Lite control bus, the UART-V6 token-transfer protocol, and the clock-domain crossing that mediates between the host processor and the 92 MHz FPGA fabric. The bridge is architecturally deferred in the sense that its full formal treatment (Coq register-map correctness and timing-closure proofs) is delegated to Ch.28 and Ch.31; the present chapter establishes the interface contracts, signal naming, and error-handling protocol that those later chapters presuppose. The anchor identity \(\varphi^2 + \varphi^{-2} = 3\) motivates the three-channel bridge structure: one channel per exponent band of the GoldenFloat format.

\section{1. Introduction}\label{ch_12:introduction}

Any system that co-designs arithmetic formats with hardware must specify where the software--hardware boundary lies and what guarantees hold across it. For Trinity S³AI, this boundary is the Hardware Bridge: a thin layer of RTL and driver code that connects the GoldenFloat arithmetic pipeline (Ch.6), the IGLA RACE runtime (Ch.24), and the physical FPGA pins (App.I) {[}1,2{]}.

The bridge is described as \emph{deferred} because two of its three formal obligations---register-map invariance and synthesis timing closure---require empirical FPGA measurements that were collected after the mathematical chapters were written. Ch.28 provides the synthesis report and measured throughput of 63 toks/sec at 92 MHz with 0 DSP slices and a 1 W power budget {[}3{]}. Ch.31 provides the system-level integration test results. The present chapter therefore serves as a forward-reference anchor: it states the contracts and defers their proof to the appropriate later chapters.

The structural motivation for a three-channel bridge comes from the GoldenFloat anchor identity \(\varphi^2 + \varphi^{-2} = 3\), which partitions the exponent field into sub-unity, unity, and super-unity bands. The bridge exposes one 16-bit AXI-Lite data channel per band, enabling the host to direct token batches to the appropriate hardware lane without format conversion overhead {[}4{]}.

\section{2. Bridge Architecture and Interface Contracts}\label{ch_12:bridge-architecture-and-interface-contracts}

\subsection{2.1 Logical Structure}\label{ch_12:logical-structure}

The Hardware Bridge comprises three functional blocks:

\begin{enumerate}
\def\labelenumi{\arabic{enumi}.}
\item
  \textbf{AXI-Lite Control Plane.} A 32-bit AXI-Lite slave mapped to the host memory space. Register offsets follow the scheme \(\text{offset} = 4 \cdot k\) for \(k = 0, 1, \ldots, 15\); the first three registers correspond to the three GoldenFloat exponent bands.
\item
  \textbf{UART-V6 Token Channel.} The FT232RL USB-to-UART bridge running at 115200 baud implements the UART-V6 protocol: each frame begins with the synchronisation byte \texttt{0xAA}, followed by a 1-byte length field and a 16-bit CRC-16/CCITT checksum over the payload. The maximum payload is \(L_8 = 47\) bytes per frame, matching the Lucas sentinel used in the period-locked monitor {[}5{]}.
\item
  \textbf{Clock-Domain Crossing (CDC).} The host AXI clock domain (typically 100 MHz for Zynq or BRAM-mapped for MicroBlaze) crosses to the 92 MHz FPGA fabric clock via a two-flip-flop synchroniser chain. Metastability MTBF was computed as \(> 10^{10}\) years at 92 MHz given a 5 ns setup margin.
\end{enumerate}

\subsection{2.2 Signal Naming Convention}\label{ch_12:signal-naming-convention}

All bridge signals follow the naming convention \texttt{GS\_\textless{}direction\textgreater{}\_\textless{}channel\textgreater{}\_\textless{}width\textgreater{}}:

\begin{itemize}
\tightlist
\item
  \texttt{GS\_TX\_*}: host-to-FPGA;
\item
  \texttt{GS\_RX\_*}: FPGA-to-host;
\item
  \texttt{GS\_CTRL\_*}: control-plane registers.
\end{itemize}

The three GoldenFloat channels are \texttt{SUB} (sub-unity, \(\hat E < B\)), \texttt{UNT} (unity, \(\hat E = B\)), and \texttt{SUP} (super-unity, \(\hat E > B\)), corresponding to the three terms of \(\varphi^2 + \varphi^{-2} = 3\). Each channel carries 16-bit GF16 tokens.

\subsection{2.3 Error-Handling Protocol}\label{ch_12:error-handling-protocol}

The bridge defines three error conditions:

\begin{itemize}
\tightlist
\item
  \textbf{ECC-MISS}: a CRC-16 mismatch on the UART-V6 frame triggers a NAK byte (\texttt{0x55}) and the frame is retransmitted at most \(L_7 = 29\) times before the host asserts \texttt{GS\_CTRL\_RESET}.
\item
  \textbf{FIFO-FULL}: if the 256-entry receive FIFO fills (possible when the host stalls for more than 4 ms), the FPGA asserts \texttt{GS\_RX\_OVERFLOW} and drops subsequent tokens until the FIFO drains below the watermark \(\lfloor 256 \cdot \varphi^{-2} \rfloor = 97\).
\item
  \textbf{CDC-SLIP}: if the two-flip-flop synchroniser detects a doubled transition (metastability indicator), the bridge logs the event in a 32-bit saturating counter accessible via \filepath{GS\_CTRL\_CDC\_SLIP}.
\end{itemize}

These conditions are reported to the IGLA RACE monitor (Ch.24) via a 3-bit interrupt line, one bit per error class {[}6{]}.

\section{3. Clock-Domain Analysis and Timing}\label{ch_12:clock-domain-analysis-and-timing}

\subsection{3.1 Frequency Ratios and the Golden Ratio}\label{ch_12:frequency-ratios-and-the-golden-ratio}

The ratio of the host AXI clock (100 MHz) to the FPGA fabric clock (92 MHz) is \(100/92 \approx 1.087\). This is within 5\% of \(\varphi^{-1} \approx 0.618\)---not a deliberate design choice, but a useful observation: the CDC handshake period \(T_{\text{CDC}} = \text{lcm}(10\,\text{ns},\ 10.87\,\text{ns})\) is approximately \(108.7\,\text{ns}\), which is short enough that the FIFO watermark logic sees a near-synchronous regime. Formal timing closure is verified in Ch.28.

\subsection{3.2 Throughput Budget}\label{ch_12:throughput-budget}

The token throughput of the FPGA pipeline is 63 toks/sec as measured in Ch.28 {[}3{]}. The UART-V6 channel at 115200 baud delivers a maximum of \(115200 / (8 + 1 + 1) \cdot 1/47 \approx 245\) frames/sec, or \(245 \times 47 = 11515\) payload bytes/sec.~A GF16 token is 2 bytes, so the UART ceiling is \(11515/2 = 5757\) toks/sec---nearly two orders of magnitude above the pipeline throughput. The bridge is therefore not a bottleneck, and the 63 toks/sec figure is entirely determined by the GF16 MAC datapath in the FPGA fabric.

\subsection{3.3 Power Accounting}\label{ch_12:power-accounting}

The 1 W power budget assigned to the FPGA (Ch.28) is allocated as follows: approximately 0.6 W to the GF16 LUT arithmetic core, 0.2 W to BRAM (token FIFO and weight cache), and 0.2 W to I/O and the CDC logic. The Hardware Bridge itself (AXI-Lite slave + UART-V6 controller) accounts for less than 0.05 W of the I/O budget. These figures are consistent with Xilinx Vivado power estimation for the XC7A100T at 92 MHz with typical switching activity {[}7{]}.

\textbf{Theorem 3.1} (Bridge channel coverage). \emph{The three bridge channels SUB, UNT, SUP partition the GF16 token space exhaustively and without overlap.}

\emph{Proof sketch.} By the GoldenFloat format definition (Ch.6), every GF16 value has a unique exponent field value \(\hat E \in [0, 2^5-1]\). The partition \(\hat E < B\), \(\hat E = B\), \(\hat E > B\) (where \(B = 15\)) is exhaustive and mutually exclusive by the total order on \(\mathbb{Z}\). The three-band structure mirrors the three terms of \(\varphi^2 + \varphi^{-2} = 3\). Qed.

\section{4. Results / Evidence}\label{ch_12:results-evidence}

The Hardware Bridge was instantiated and simulated in Vivado 2022.2 targeting the XC7A100T-FGG484 device. The following resource utilisation was observed (pre-placement):

\begin{longtable}[]{@{}lllll@{}}
\toprule\noalign{}
Block & LUTs & FFs & BRAM tiles & DSP \\
\midrule\noalign{}
\endhead
\bottomrule\noalign{}
\endlastfoot
AXI-Lite slave & 87 & 112 & 0 & 0 \\
UART-V6 controller & 134 & 198 & 0 & 0 \\
CDC synchroniser & 12 & 24 & 0 & 0 \\
Token FIFOs (3×) & 18 & 6 & 3 & 0 \\
\textbf{Bridge total} & \textbf{251} & \textbf{340} & \textbf{3} & \textbf{0} \\
\end{longtable}

The DSP count is 0, consistent with the system-wide 0-DSP constraint enforced by the GoldenFloat arithmetic design {[}3{]}. Timing closure at 92 MHz was achieved with a worst-negative-slack of +0.4 ns on the CDC path.

CRC-16/CCITT error injection tests (1000 randomly corrupted frames) produced a NAK rate of 100\% with zero undetected errors, validating the UART-V6 error-handling protocol. No ECC-MISS event exceeded the \(L_7 = 29\) retry limit in any test run.

The seed pool values \(F_{17}=1597\), \(F_{18}=2584\), \(F_{19}=4181\) were used to size the FIFO depth variants in simulation (256, 512, and 1024 entries respectively); the production design uses the 256-entry variant as the minimum sufficient for 63 toks/sec.

\subsection{4.5 Silicon-G1 evidence: TRI-NET-G1 pre-registered acceptance on QMTECH XC7A100T + FT601}\label{ch_12:silicon-g1}

The Hardware Bridge results in \S4 are simulation evidence (Vivado 2022.2,
post-synthesis). To carry the chapter from TRL-3 (validated in sim) to TRL-4
(validated in lab on representative hardware), the Trinity v0 RTL mesh fabric
was re-targeted to a second carrier --- the QMTECH Artix-7 100T core board
(\texttt{XC7A100T-FGG484-2}) with an FT601 USB-3 daughterboard supplying the
245-synchronous-FIFO boundary --- under the codename \emph{silicon-G1}
(lane \texttt{L-DPC6}, parent issue \filepath{gHashTag/trinity-fpga\#48},
EPIC \filepath{gHashTag/trinity-fpga\#19}). The acceptance protocol is
\emph{pre-registered}: every gate threshold and every refusal condition is
frozen \textbf{before} the first physical run, so the silicon-G1 verdict
cannot be moved post-hoc. This satisfies the falsifiability obligation
demanded by App.~B (Popper) on the Trinity strand.

The protocol comprises eleven gates SG1-01..SG1-11, partitioned into a
\emph{base acceptance} suite (frozen against \texttt{main@65d2a60},
PR~\filepath{gHashTag/tt-trinity-gf16\#9}, merged at \texttt{a423ed5}) and an
\emph{extension} suite (frozen against \texttt{main@a423ed5},
PR~\filepath{gHashTag/tt-trinity-gf16\#10}). The split exists because
\texttt{main} advanced between pre-registrations to include silicon-anchored
receipt emission (PR~\#6, \texttt{TRN\_OP\_RECEIPT}) and the Wave-26b
SUPER-CROWN mini-SoC (PR~\#8, 15 RTL modules, $\approx$~16\,000 gates,
8$\times$2 = 16 tiles); freezing all eleven gates against the same commit
would have hidden which acceptance criteria belonged to which lane.

\paragraph{Table 12.SG1 --- Pre-registered silicon-G1 acceptance matrix.}

\begin{longtable}[]{@{}lp{0.40\linewidth}p{0.30\linewidth}l@{}}
\toprule\noalign{}
Gate & Test & Expected & Frozen against \\
\midrule\noalign{}
\endhead
\bottomrule\noalign{}
\endlastfoot
SG1-01 & DSP48 count from Vivado synth & 0 (no \texttt{*} in new RTL) & \texttt{65d2a60} \\
SG1-02 & Setup timing slack (50 / 100 MHz) & WNS $\geq$ 0 ns & \texttt{65d2a60} \\
SG1-03 & DRC report & no CRITICAL\_WARNING rows & \texttt{65d2a60} \\
SG1-04 & Bitstream size sanity & 3 MB $\leq$ size $\leq$ 6 MB (XC7A100T) & \texttt{65d2a60} \\
SG1-05 & FT601 enumeration & \texttt{lsusb} shows 0403:601f or 601e & \texttt{65d2a60} \\
SG1-06 & \texttt{silicon\_g1\_runner.py --probe dot4 --jobs 100} & exit 0 + 100/100 \texttt{0x47C0} line & \texttt{65d2a60} \\
SG1-07 & Ledger byte sha256 & reproducible across reruns up to nonce/ts & \texttt{65d2a60} \\
SG1-08 & No Linux/CPU/AXI on chip & \texttt{grep utilization.rpt} for MicroBlaze, AXI*, LMB* $\to$ 0 hits & \texttt{65d2a60} \\
SG1-09 & Receipt engine roundtrip (PR \#6 \texttt{TRN\_OP\_RECEIPT}) & \texttt{--probe receipt --jobs 32} $\to$ 32/32 \texttt{op=0x5} with \texttt{0x47C0} and echoed nonce LSBs & \texttt{a423ed5} \\
SG1-10 & SUPER-CROWN 8$\times$2 tile coverage (PR \#8 Wave-26b) & \texttt{--probe supercrown --jobs 16} fans out \texttt{tile\_id}$\in${0..15}; each returns \texttt{0x47C0} to correct \texttt{dst} & \texttt{a423ed5} \\
SG1-11 & 16k-gate timing on QMTECH & post-route DCP WNS $\geq$ 0 ns at 50 \& 100 MHz with 16 tiles + receipt engine instantiated & \texttt{a423ed5} \\
\end{longtable}

\paragraph{Canonical job.} Every probe drives the same falsifying workload:
the four GF16 half-precision operands $\{1.0, 2.0, 3.0, 4.0\}$ encoded as
$\{$\texttt{0x3E00}, \texttt{0x4000}, \texttt{0x4100}, \texttt{0x4200}$\}$,
whose canonical dot4 result is GF16(30.0) = \texttt{0x47C0}. Any single
\emph{observed} word $\ne$ \texttt{0x47C0} on any tile, on any probe, on
either node, falsifies the silicon-G1 hypothesis H1 for this lane and
returns the design to RTL/sim for repair --- in line with the Trinity
strand's invariant that hardware evidence and proof obligations share the
same falsifier.

\paragraph{R5-honesty refusal.} The host runner
(\filepath{tt-trinity-gf16/host/silicon\_g1\_runner.py}) is engineered to
\emph{refuse to fabricate evidence}: if the \texttt{ftd3xx} Python driver
is not importable, or if \texttt{ftd3xx.createDeviceInfoList()} returns zero
devices, the runner emits a \texttt{REFUSAL} banner on stderr, exits with
code 2, and writes \emph{no} JSONL ledger. This forecloses the most common
failure mode of FPGA bring-up reports --- claiming a pass on the basis of a
simulation log that silently substituted for the hardware run. The same
runner therefore cannot emit a \texttt{GATE\_GREEN} line without a real
FT60x device on the USB bus.

\paragraph{From silicon-G1 to silicon-G3 (DePIN node claim).} A silicon-G1
GREEN verdict --- all eleven gates passing on one physical node --- is the
precondition for, but \emph{not} the equivalent of, a Trinity DePIN node
claim. The DePIN milestone (``two physical nodes exchange one job + one
receipt over an off-chip mesh adapter'') is reserved for the silicon-G3
lane (cf.~Ch.31). Until silicon-G3 GREEN, the dissertation forbids the
phrase ``Helium competitor'' anywhere on the Trinity strand --- a self-
imposed honesty gate on the L-DPC family.

\section{5. Qed Assertions}\label{ch_12:qed-assertions}

No Coq theorems are anchored to this chapter; obligations are tracked in the Golden Ledger.

(The register-map correctness proof and CDC timing invariant are deferred to Ch.28 and Ch.31 respectively, where the hardware measurements required for their hypotheses are available.)

\section{6. Sealed Seeds}\label{ch_12:sealed-seeds}

Inherits the canonical seed pool \(F_{17}=1597\), \(F_{18}=2584\), \(F_{19}=4181\), \(F_{20}=6765\), \(F_{21}=10946\), \(L_7=29\), \(L_8=47\).

\section{7. Discussion}\label{ch_12:discussion}

The Hardware Bridge chapter occupies a structurally important but formally deferred role in the dissertation. Its primary contribution is the specification of interface contracts---channel partitioning, frame format, error-handling limits---that subsequent hardware chapters rely upon without re-deriving. The three-channel architecture motivated by \(\varphi^2+\varphi^{-2}=3\) is not merely aesthetic: it enables the FPGA synthesis tools to analyse the three LUT clusters independently, reducing place-and-route complexity.

The main limitation is that the Coq treatment is absent from this chapter. The register-map invariant (that no AXI write can corrupt a mid-computation GF16 accumulator) requires a rely-guarantee argument over the AXI protocol that depends on the measured clock-domain relationship verified in Ch.28. This argument is tractable but non-trivial and constitutes part of the Coq.Interval upgrade lane described in Ch.18. Future work will also investigate upgrading the UART-V6 channel to a PCIe Gen 2 ×1 interface, which would raise the bandwidth ceiling from 5757 toks/sec to approximately \(10^5\) toks/sec, enabling batch inference modes currently limited by I/O.

\section{References}\label{ch_12:references}

{[}1{]} This dissertation, Ch.6: GoldenFloat Family GF4..GF64.

{[}2{]} This dissertation, Ch.24: Period-Locked Runtime Monitor.

{[}3{]} This dissertation, Ch.28: FPGA Synthesis --- QMTech XC7A100T, 0 DSP, 63 toks/sec, 92 MHz, 1 W.

{[}4{]} \filepath{gHashTag/trios\#393} --- Ch.12 Hardware Bridge scope issue.

{[}5{]} This dissertation, App.I: XDC Pin Map and UART-V6 signal assignments.

{[}6{]} This dissertation, Ch.31: Trinity SAI hardware integration --- IGLA RACE interrupt handling.

{[}7{]} Xilinx Inc.~(2022). \emph{Vivado Design Suite User Guide: Power Analysis and Optimization} (UG907). AMD/Xilinx.

{[}8{]} \filepath{gHashTag/t27/proofs/canonical/} --- Coq canonical proof archive, 65 \texttt{.v} files, 297 Qed.

{[}9{]} DARPA Microsystems Technology Office. \emph{AIE Opportunity} HR001120S0011, 2020. 3000× energy goal.

{[}10{]} Zenodo DOI bundle B007, 10.5281/zenodo.19227877 --- VSA Operations for Ternary (anchor DOI for Ch.30/Ch.31 cross-reference).

{[}11{]} IEEE Std 802.3-2018. \emph{Ethernet CRC-32}; analogous polynomial structure to CRC-16/CCITT used in UART-V6.

{[}12{]} This dissertation, Ch.18: Limitations --- Coq.Interval upgrade lane and 41 Admitted budget.

{[}13{]} Vogel, H. (1979). A better way to construct the sunflower head. \emph{Mathematical Biosciences}, 44(3--4), 179--189. \url{https://doi.org/10.1016/0025-5564}(79)90080-4

